% ============================================================
% Acknowledgments (الشكر والعرفان) - Arabic Format
% ============================================================
\thispagestyle{plain}
\vspace*{0.5cm}

\begin{center}
\begin{Arabic}
{\fontsize{18}{22}\selectfont\bfseries الشكر والعرفان}
\end{Arabic}
\end{center}
\vspace{0.6cm}

\begin{Arabic}
\fontsize{12}{22}\selectfont
\setstretch{1.6}

الحمد لله رب العالمين، حمداً يليق بجلال وجهه وعظيم سلطانه، على ما أنعم به علينا من توفيق وهداية، وما وهبنا من قدرة وصبر حتى وصلنا إلى إتمام هذا البحث وتكليل مسيرتنا العلمية بهذا الإنجاز.

وعملاً بالهدي النبوي الشريف: «مَنْ لا يَشْكُرُ النَّاسَ لا يَشْكُرُ اللَّهَ»، يُسعدنا ويشرّفنا أن نتقدم ببالغ الشكر وعظيم الامتنان إلى مشرف مشروعنا الفاضل \textbf{الدكتور نصر قايد علي العودي}، الذي كان لنصائحه القيمة وتوجيهاته الأكاديمية السديدة ومتابعته الدؤوبة أبلغ الأثر في تذليل الصعاب وتوجيه مسار هذا العمل الهندسي نحو أهدافه المرجوة.

كما نتقدم بخالص الشكر والتقدير إلى جميع أساتذتنا الأفاضل في \textbf{قسم الهندسة الحيوية الطبية} و\textbf{كلية الهندسة} في جامعة العلوم والتكنولوجيا، الذين لم يدخروا جهداً في تزويدنا بالعلوم والمعارف وإرساء القواعد العلمية والهندسية المتينة طوال سنوات دراستنا الجامعية.

ونتوجه بالشكر إلى كل من قدم لنا يد العون والمساعدة، أو أسدى إلينا نصيحة تقنية، أو دعمنا بكلمة مشجعة كان لها أثر طيب في إنجاح هذا العمل.

نسأل الله العلي القدير أن يبارك في هذا الجهد، وأن يجعله نافعاً وخالصاً لوجهه الكريم، وأن يجزيهم جميعاً عنا خير الجزاء.

\end{Arabic}

\normalsize
\clearpage
